\section{Background}

\subsection{Syntax and Semantics of \ltlf}
\label{sec:ltlf}

Given a finite set $\calP$ of atomic propositions, an \ltl formula $\varphi$ has the form:
%\begin{equation}
\[
    \varphi ::= \top \mid
    \bot
    \mid 
    a
    \mid
    \lnot\varphi
    \mid
    \varphi_1\land\varphi_2
    \mid
    X\varphi
    \mid
    \varphi_1 U\varphi_2,
\]
%\end{equation}
where $a \in \calP$. We use $\top$ and $\bot$ to denote $True$ and $False$, respectively. $X$ (Next) and $U$ (Until) are temporal operators. Other temporal operators can be derived as follows: $R$ (Release) is defined as $\varphi_1 R \varphi_2 \equiv \neg(\neg\varphi_1 U \neg\varphi_2 )$, $G$ (globally) is defined as $G\varphi \equiv \bot R\varphi$, and $F$ (eventually) is defined as $F \varphi \equiv \top U \varphi$.
%
\ltlf is a finite-trace language derived from \ltl and share the same grammar, with the difference that $X$ is understood as Strong Next and require that the next time step in a given trace at a given position exists. In contrast, Weak Next is denoted by $W$ and does not require the existence of the next time step, being interpret as $\top$ if evaluated at the last time step of the trace. It can be expressed in terms of Strong Next by $W \varphi \equiv \neg X\neg\varphi$. 

A trace $\boldsymbol{\sigma} = (\sigma_{0}, \sigma_{2}, \dots, \sigma_{T})$ is a sequence of propositional assignments to the propositions in $\calP$, where $\sigma_{t} \in 2^\calP$ is the set of all and only propositions that are true at instant $t$. We use $|\boldsymbol{\sigma}|$ to denote the length of the trace, and $\boldsymbol{\sigma}[0,i]$ to denote a proper prefix of $\boldsymbol{\sigma}$ from time step $0$ to time step $i$.
%Since every trace is finite, $|\boldsymbol{\sigma}| < \infty$.
Note that, if the propositional symbols in $\calP$ are all \textit{mutually exclusive}, i.e., the domain produces exactly one symbol true at each step, then we have $\sigma_{t} \in \calP$.
%
Let $last = |\boldsymbol{\sigma}| - 1$. We inductively define when
an \ltlf formula $\varphi$ is true at an instant $i$ (for $0 \leq i \leq last$), written $\boldsymbol{\sigma},i \models \varphi$, as follows:
\begin{itemize}
    \item $\boldsymbol{\sigma}, i \models a$ iff $a \in \boldsymbol{\sigma}_i$, for $a \in \calP$;
    \item $\boldsymbol{\sigma}, i \models \neg \varphi$ iff $\boldsymbol{\sigma}, i \not\models \varphi$;
    \item $\boldsymbol{\sigma}, i \models \varphi_1 \land \varphi_2$ iff $\boldsymbol{\sigma}, i \models \varphi_1$ and $\boldsymbol{\sigma}, i \models \varphi_2$;
    \item $\boldsymbol{\sigma}, i \models X\varphi$ iff $i < last$ and  $\boldsymbol{\sigma}, i + 1 \models \varphi$;
    \item $\boldsymbol{\sigma}, i \models \varphi_1 \muntil \varphi_2$ iff for some $j$ such that $i \leq j \leq last$, we have $\boldsymbol{\sigma}, j \models \varphi_2$, and for all $k$ such that $i \leq k < j$, we have $\boldsymbol{\sigma}, k \models \varphi_1$.
\end{itemize}
%
\noindent
We say that $\boldsymbol{\sigma}$ satisfies $\phi$, written $\boldsymbol{\sigma} \models \varphi$, if $\boldsymbol{\sigma}, 0 \models \varphi$. A formula $\varphi$ is satisfiable if it is true wrt some trace $\boldsymbol{\sigma}$, and is valid, if it is true wrt every trace $\boldsymbol{\sigma}$.
%
On the empty trace $\varepsilon$, Boolean connectives behave as usual, atoms are false, $U$ and $X$ are false, and $R$ and $W$ are true. Note that $F\top$ holds exactly on non-empty suffixes of a trace, and $G\bot$ exactly on the empty one.


In the following, we will use $sub(\varphi)$ to denote the set of proper subformulas of $\varphi$, and we will refer to it as the syntactic closure of $\varphi$. Given a formula $\varphi$ and a propositional assignment $\sigma_t$ at time $t$, we are also interested in computing the residual formula, i.e the formula that has to be satisfied at time $t+1$. As explained in details in Sec. \ref{sec:rulerunner-progression}, a residual formula can be computed using formula progression \cite{}, which returns a positive Boolean formula on $sub(\varphi) \cup sub(F \top) \cup sub(G \bot)$. We will use $cl(\varphi)$ to denote the set of residual formulas, and we will refer to it as the residual closure of $\varphi$ .

\subsection{Symbolic Monitoring}
\label{sec:symbolic}

\paragraph{Deterministic Finite Automata.}
Any \ltlf formula $\varphi$ can be translated into a Deterministic Finite Automaton (DFA) \cite{LTLf} $A_\varphi = (2^\calP, Q, q_0, \delta, F)$, where $2^\calP$ is the automaton alphabet, $Q$ is the finite set of states, $q_0 \in Q$ is the initial state, $\delta: Q \times 2^\calP \rightarrow Q$ is the transition function, and $F \subseteq Q$ is the set of final states. 
We also define the extended transition function over traces $\delta^*: Q \times (2^\calP)^* \rightarrow Q$ recursively as:
\begin{equation}
%\[
\begin{array}{l}
\delta^*(q,\epsilon) = q \\
\delta^*(q, \sigma+\boldsymbol{x}) = \delta^*(\delta(q,\sigma) , \boldsymbol{x}),
\end{array}
%\]
\end{equation}
where $\sigma \in 2^\calP$ is a symbol and $\boldsymbol{x} \in (2^\calP)^*$ is a trace.  
The automaton accepts the trace $\boldsymbol{\sigma}$ iff $\delta^*(q_0, \boldsymbol{\sigma}) \in F$, and in that case we say that $\boldsymbol{\sigma}$ belongs to the language of the automaton, denoted as $\calL(A_\varphi)$.  
We have that $\varphi$ and $A_\varphi$ are equivalent because, for any trace $\boldsymbol{\sigma} \in (2^\calP)^*$:
\begin{equation}
%\[
\boldsymbol{\sigma}\in \calL(A_\varphi) \iff\boldsymbol{\sigma}\vDash\varphi.
%\]
\end{equation}

We can rely on the use of DFAs for several reasoning tasks, like model checking, planning and synthesis. In particular, monitoring formula $\varphi$ becomes a deterministic walk over the automaton $A_\varphi$: the monitor keeps a single current state $q$, initialized to $q_0$, and on observing the time point $\sigma_t \in 2^\calP$ it advances $q \leftarrow \delta(q, \sigma_t)$. Because the automaton is fixed and deterministic, each step costs a single transition lookup, independent of the trace seen so far and of the size of $\varphi$. In the following, DFA-based monitoring will serve as the exact baseline against which the neuro-symbolic paradigms will be measured.

%as implemented by \texttt{ltlf2dfa} \cite{fuggitti-ltlf2dfa}.
%the automaton encodes the specification exactly and admits no learning or adaptation.

\paragraph{Verdicts and early termination.}
To address online monitoring, we can resort to a three-valued runtime semantics \cite{monitorability_LTL}: at each step the monitor emits \textsc{satisfy}, \textsc{violate}, or \textsc{undecided}, the last meaning that the observed prefix can still be extended into both an accepting and a non-accepting continuation. A definite verdict can be issued as soon as the outcome is settled regardless of the future, which corresponds to two structural properties of the current state, both precomputed once by backward reachability over $A_\varphi$. A state is a \emph{trap} if no accepting state is reachable from it: reaching a trap means $\varphi$ is permanently violated, so the monitor halts with \textsc{violate}. Dually, a state is an \emph{accepting sink} if every state reachable from it (including itself) is accepting: the monitor halts with \textsc{satisfy}. With trap and sink sets materialized at compile time, the per-step verdict reduces to two set-membership tests on top of the transition lookup. As long as the current state is neither, the verdict is \textsc{undecided} and monitoring continues.
Some specifications never reach a trap or an accepting sink, so \textsc{step} returns \textsc{undecided} for the entire trace. An example is the response pattern $G(a \rightarrow Fb)$, which can always be satisfied by a later $b$ and always be rescued from violation, hence has neither an accepting sink nor a trap. For these the verdict is emitted at the end of the trace, and the monitor reports it by testing acceptance of the state reached, i.e.\ \textsc{satisfy} iff $\delta^*(q_0, \boldsymbol{\sigma}) \in F$ and \textsc{violate} otherwise. Note that, in online monitoring, we require the monitor to be prefix-sound, i.e. if it
outputs \textsc{satisfy} (resp.\ \textsc{violate}) after reading the prefix $\boldsymbol{\sigma}[0,t]$, then
$\boldsymbol{\sigma}'\models\theta$ (resp.\ $\boldsymbol{\sigma}'\not\models\theta$)
for every trace $\boldsymbol{\sigma}'$ sharing the same prefix.
This is a stronger property than being able to emit the correct verdict at the end of the trace.